\cotefile{156-3}
\scanpage{26}\mspage{25}
% L1–4
\unsure{pose}
\[
\partial T=\text{bord absolu de }T
=\left\{\deleted{t\in T\mid t\text{ est }}\begin{aligned}
&\text{un plus petit}\\[-2pt]
&\text{ou plus grand}\\[-2pt]
&\text{élément de }T
\end{aligned}\right\}
\]
% L5–6
\deleted{Si $T$ est réduit à un pt,}

On a donc
% L7–10
\[
\begin{aligned}
\partial\{a\}&=\{a\}\\
\partial[\unsure{o|a|0},b]&=\{a,b\} &&\text{si }a\leq b\\
\partial L_{\geq a}&=\partial L_{\leq a}=\{a\} &&\forall\,a\in L\\
\deleted{\unsure{\partial L=\varnothing}}\quad\partial L&=\partial L\quad\text{(sic)}
\end{aligned}
\]
\margin{(sans exclure le cas $a\in\partial L$)}
% L11–12
Notons que
\[
\unsure{T'|\dot T}\cap T\subset\partial T\subset\deleted{\partial L}(T'\cap T)\cup\partial L
\]
% L13–15
i.e tout $t\in T$ qui est pt frontière, est dans le bord \inserted{de $T$}, et l'inverse est vrai si $t\notin\partial L$. Plus précisément
% L16–18
\[
\begin{aligned}
&\deleted{\text{Si }\unsure{\operatorname{Card}}\,T=1_{\unsure{\mathrm{elt}|\mathrm{él}}},\quad\partial T=\dot T=T}\\
&\deleted{\text{Sinon}}\quad
\left\{\begin{aligned}
&\deleted{T'\cap T=\partial T\setminus(\partial T\cap\partial L)}\quad\deleted{\text{\unsure{i.e.}}}\\
&\deleted{\partial T=(T'\cap T)\cup T\cap\partial L}
\end{aligned}\right.
\end{aligned}
\]
% L19
\[
\boxed{\partial T=(T'\cap T)\mathbin{\unsure{\underline{\cup}|\cup}}(T\cap\partial L)}
\]
% L20–22
\unsure{Tronçons|Tronçon} : C'est une partie \deleted{\unsure{de la réunion}} \deleted{de la forme} \deleted{\unsure{convexe}} de la forme $T'=T\setminus\alpha$, avec $\alpha\subset\partial T$.

\deleted{\unsure{Tronçon ouvert|Tronçons ouverts} :} \unsure{c'est|i.e.} un ens de la forme
% L23–26
\[
\left\{\begin{aligned}
&]a,b],\;[a,b[,\;]a,b[,\quad(a<b)\\
&L_{<a},\ \deleted{\unsure{a\in\partial L})}\quad(a\neq\unsure{\omega|\infty})\\
&L_{>a}\hspace{7em}(a\neq\alpha)\\
&L
\end{aligned}\right.
\]
% slanted
% L lower left, oblique
\margin{\uline{NB} Cette écriture est \uline{unique} \unsure{on|On} pose $T=\widetilde{T'}$ (c'est une \unsure{sorte} d'adhérence \unsure{de $T$|de $T'$} -- \unsure{c'est} le \unsure{sous-tronçon|seul tronçon} le plus petit fermé \unsure{qui contient} $T'$.}
\scanpage{27}\mspage{26}
% L1–5
Les tronçons $]a,b[,\ L_{<a},\ L_{>a},\ \inserted{L}$ sont dits ouverts \deleted{(\uline{NB} ils sont} \deleted{(ce sont bien ceux} \inserted{(ce sont bien ceux qui sont} ouverts pour la top. can.), \deleted{qui sont ouverts} et $]a,b]$ et $[b,\unsure{o|a|0}[$ sont dits \uline{demi}-\uline{ouverts}.

% L6–11
\uline{Attention} \deleted{Le tronçon} Si $a$ plus petit \unsure{élém|élément}, alors $[a,b[$, \deleted{est ouvert} \inserted{qui est} demi-ouvert pour la terminologie \unsure{précédente}, est \unsure{aussi} ouvert pour la top. canonique, et il est égal à $L_{<b}$. Kif kif \unsure{dualement|autrement}. $L$ est à la fois ouvert et fermé.

% L12–14
Un \uline{tronçon} \uline{générali\deleted{\unsure{sé}}\unsure{sé}} \inserted{ou \unsure{ps.tronçon}} est une partie de $L$ qui est soit \unsure{réduit|réduite} à un pt, soit un tronçon. Pour un tronçon de la forme
% L15–16
\[
\deleted{\unsure{T}}\ T'=T\setminus\alpha\qquad(T\text{ tronçon fermé},\ \alpha\in\partial T
\]
je pose
\[
\partial T'=\partial T
\]
% L17–18
et \unsure{on aura|l'on aura} \unsure{encore|que}
\[
\partial T'=T'\mathbin{\unsure{\dot\cup|\dot\cap|\cup}}\partial L\deleted{\unsure{.}}
\]
\margin{Attention, cette fois ce n'est pas le $\partial$ absolu de $T'$ pour la structure d'ordre induite. \unsure{\ldots}}
% slanted
% L lower left, oblique
\margin{\uline{NB} $\partial T'\neq\varnothing$ sauf \unsure{qd|q} si $T'=L$, $L$ \unsure{une|un} \deleted{\unsure{``droite ouverte''}} \deleted{\unsure{ligne}} \unsure{tronçon-interminable}\ldots\ Les pts frontière de $T'$ sont \deleted{ses} bornes inf. et sup. \uline{dans $T$} (si elles existent) : borne sup (ssi $T'$ est majoré) borne inf (ssi $T'$ est minoré) égales \unsure{ds} le cas $T$ \unsure{réduit} à un pt\ldots}

% L19–22
\deleted{\unsure{Soient} deux} Deux \unsure{ps.tronçons} \inserted{$T,T'$} sont dits \uline{en pos. relative modérée} (\unsure{abrég}, $(T,T')$ \deleted{\unsure{en p.}} modéré ou $\{T\unsure{\mid|,}T'\}$ modéré) si $\partial T$ et $\partial T'$ le sont, i.e $\partial T\cup\partial T'$ est tot. ordonné.
\scanpage{28}\mspage{27}
% L1–3
\noindent\llap{\rule[-1.05\baselineskip]{0.4pt}{1.85\baselineskip}\hspace{0.65em}}\uline{Prop} Si $T,T'$ sont des ps. tr. en position rel modérée, alors $T\cap T'$ est \inserted{vide ou est} un ps. tr, et on a
\[
\partial(T\cap T')\subset\partial T\cup\partial T'
\]

% L4–6
Immédiat. On va introduire une relation d'ordre entre ps-tronçons \inserted{(\unsure{pas|pour} l'inclusion!)}. On dit que
\[
T\unsure{<|\prec}T'
\]
% L7–8
si \deleted{a) $T$ est en pos. rel modérée avec $T'$,}
\[
\deleted{\text{b)}}\quad\deleted{\unsure{\forall}}\quad t\in T,\ t'\in T'\implies t<t'
\]
\margin{\leavevmode\llap{\rule[-4\baselineskip]{0.35pt}{5\baselineskip}\hspace{0.4em}}\deleted{\unsure{Si les}} Ça implique $T$ et $T'$ en pos. rel. mod et $T\cap T'=\varnothing$}
% L9
C'est \unsure{évident|évidemt} transitif.

% L10–15
Soit $T<T'$. Comme $T$ est majoré, il s'ensuit que \deleted{$\partial T\neq\varnothing$} il a un plus grand \unsure{élém|élément} \deleted{\unsure{b}} $a$, \deleted{tel que} \inserted{et alors} $T\subset L_{\leq a}$. \deleted{\unsure{Comme}} De $=$, $\partial T'$ a un plus petit élément $b$, \unsure{tel} que $T'\subset L_{\geq b}$, \deleted{\unsure{et on a}} \deleted{\unsure{alors}} et on voit de suite que $a\leq b$.

% L16 and note below brace
\uline{Position relative \deleted{\unsure{impropre}} \inserted{de raccordement}} : $a=b$, et $\underbrace{a\in T\text{ ou }b\in T'}$%
\margin{à l'exclusion l'un de l'autre, puisque $T\cap T'=\varnothing$}

% L17–20
\deleted{\unsure{Pour|Par} la suite, on l'\uline{exclut} \unsure{donc} :}

On renforce la relation $<$ en $\ll$
\[
\begin{aligned}
T\ll T' &\Longleftrightarrow T<T',\quad\text{et position relative}\\
&\hspace{12em}\text{propre}\\
&\Longleftrightarrow \exists\,x\in L\text{ tel que }T\subset L_{\unsure{<|\leq}\unsure{a|x}},\ T'\subset L_{>x}
\end{aligned}
\]
\scanpage{29}\mspage{28}
% L1
\uline{psTronçon} \uline{interstitiel}

% L2–3
\uline{Prop} Soit $T,T'$ \deleted{ps. tr. tels que $T<T'$}, \deleted{$a=\operatorname{Sup}T,\ b=\operatorname{Inf}T'$}, Conditions équivalentes \deleted{si $a\in{}$} \inserted{(si $a=\operatorname{Sup}T,\ b=\operatorname{Inf}T'$)}

% L4–5
a) $T\unsure{<|\prec}T'$ \inserted{et $T,T'$} ne se raccordent pas, i.e.
\[
(a\in T\text{ ou }b\in T')\implies a\neq b
\]
% L6
\deleted{b) $\exists\,x\in L$}\quad b) $\exists\,x\in L$ tel que $T\subset L_{<\unsure{x|a}},\ T'\subset L_{>\unsure{x|a}}$
% slanted
\margin{\uline{NB} b) implique déjà que $T<T'$}

% L7–9
\deleted{\unsure{b|6}}\unsure{c)|b)} $\exists$ ps. tr. $\deleted{\unsure{I}}S$ tel que
\[
T<\deleted{I}S<T'
\]
et tel que $T$ et $\deleted{I}S$, $\deleted{I}S$ et $T'$ \uline{se raccordent}.

% L10–12
De plus, ce\deleted{\unsure{t}} $\deleted{I}S$ est unique, et $\partial S=\{a,b\}$. $S$ est formé des $x\in L$ tels que
\[
T\subset L_{<\unsure{a|x}},\ T'\subset L_{>x}\qquad\text{(cf b) )}
\]

% L13–14
On appelle $S$ le ps. tr. interstitiel entre $T$ et $T'$

% L15–19
\uline{Cor} 1 Soit $T$ un tronçon, et supposons que $T$ soit minoré \inserted{\unsure{strictemt|strictement}}. Alors l'ens des \inserted{minorants} \inserted{stricts, i.e. des} $x\in L$ tels que $\unsure{\overline T|T}\subset L_{>x}$ (ens des minorants stricts) est un \deleted{tronçon} ps. tr. $S$, \deleted{\unsure{que}} \inserted{qui satisfait}
% L20
\begingroup
\makeatletter
\tagsleft@true
\makeatother
\[
S<T,\ S\text{ et }T\text{ se raccordent},\ S\text{ sans minorant}\tag*{\(\ast\))}
\]
\endgroup
% L21
et ces conditions \deleted{\unsure{satisf}} caractérisent $S$.

% L22–24
\deleted{\uline{Cor} 2 \unsure{Énoncé|Enoncé} \unsure{de|d}} On appelle $S$ le \uline{pstronçon interstitiel gauche} de \deleted{\unsure{$S$}}$T$, On définit de $=$ (si $T$ majoré) le pstr. interst. droit. de $T$.
\scanpage{30}\mspage{29}
% L1–3
\deleted{\uline{Définition}}

Soit $\operatorname{Tr}(L)$ l'ens \deleted{ordonné} des ps. tronçons de $\unsure{\mathcal L|L}$, ordonné par $\ll$.

% L4–6
Pour toute partie finie tot. ordonnée $\unsure{\Phi|\varphi}$ de $\operatorname{Tr}(L)$, considérons
\[
\deleted{U}X_{\Phi}=\bigcup_{i\in\Phi}T_i\qquad\text{(NB réunion disjointe)}
\]
% L7–12
% slanted
\margin{\uline{Prop}}
Alors l'application
\diagram{\begin{tikzcd}[ampersand replacement=\&,column sep=large,row sep=1.3ex]
\Phi \arrow[r,mapsto] \& \deleted{U}X_{\Phi}\\
\underset{\substack{\text{drapeaux de}\\\operatorname{Tr}(L),\ \text{y inclus}\\\text{\unsure{par|pour|pav|pau}}}}
{\underline{{\unsure{\operatorname{Drap}}}_{\ll}^{\unsure{*|b|\varnothing}}(\operatorname{Tr}(L))}}
\arrow[r] \& \unsure{\mathcal P|\mathscr P}(L)
\end{tikzcd}}
est \uline{injective}.

% L13–15
Les ens de la forme $U_{\Phi}$ \unsure{s'appellent} \uline{modérés}. Pour un tel
\[
\deleted{U}X_{\Phi}=\bigcup T_i
\]
on pose
% L16–17
\[
\unsure{\operatorname{Fr}}(\deleted{U}X_{\Phi})=
\underbrace{\left(\bigcup_i\partial T_i\right)}_{\substack{\deleted{\text{\unsure{est}}}\\\partial U,\ \text{drapeau de } L}}\setminus\varepsilon
\]
% slanted
\margin{``frontière'' relative dans $L$ (mais pas pour la topologie \unsure{de|sur} \unsure{m.|sur} $L$\ldots)}
% L18–24
où
\[
\begin{aligned}
\varepsilon={}&\bigl\{t\in\partial L\cap U\mid
  \deleted{\unsure{t\in U}}\ \{t\}\notin\Phi\quad\text{i.e. }\deleted{\unsure{\Phi|\notin}}\\
&\hspace{6em}\{t\}\text{ n'est pas parmi les }T_i,\\
&\hspace{6em}\deleted{\text{i.e. }\exists\,x\in L,\ \text{\unsure{tq} } x<{}}\\
&\hspace{6em}\deleted{\text{i.e. il y}}\quad\text{i.e }U\text{ est un}\\
&\hspace{8em}\text{voisinage de }t\text{ dans } L\\
&\hspace{6em}\text{i.e. }\{t\}\text{ pas point isolé}\\
&\hspace{9em}\text{de }U\ \ldots\bigr\}
\end{aligned}
\]
% slanted
% L lower left note
\margin{Les \unsure{$T_i$|$X_i$} s'appellent les ``composantes (connexes!) de \deleted{\unsure{$U$}}$X$ (mais pas pour la top canonique) \unsure{l'on|l'ens} \unsure{va|en} \unsure{aussi|sera} \unsure{les|se} \unsure{noter|noté} parfois $\unsure{\Pi|\pi}_{\unsure{0|o}}(X)$}
\scanpage{31}
\mspage{30}

% L1–5
\deleted{\unsure{Ainsi|Disons}} Soit $\deleted{U}X=\deleted{U}X_{\unsure{\Phi|\varphi}}$ une partie modérée. Un $x\in L$ est dit en pos. rel. modérée avec $\deleted{U}X$, s'il l'est avec les $T_i$ (\inserted{\uline{tronçons}} \uline{composants} de $\deleted{U}X$) i.e. s'il l'est avec $\underset{\partial U}{\underline{\bigcup\partial T_i}}$, ou \unsure{encore} avec $\operatorname{Fr}(\deleted{U}X)$.
% Lleft margin, alongside lines 1–13
% slanted
\margin{Plus \unsure{généralemt|généralement}, Si $\deleted{U}X\unsure{,|\cap}\deleted{\unsure{V|U'}}Y$ sont \unsure{modérées|modérés}, on dit que $\deleted{U}X,Y$ sont en pos. rel. modérée si \deleted{$\partial U$ et $U'$}, $\partial\deleted{U}X,\partial\deleted{V}Y$ le sont ou encore $\operatorname{Fr}\deleted{U}X,\operatorname{Fr}Y$ le sont. Ex: si $\deleted{U}X\subset Y$, alors $(\deleted{U}X,\deleted{V}Y)$ \unsure{modérée|modéré}.}
% L5–7
Par exemple \deleted{Soit} $x\in\deleted{U}X$, ou $x\in\deleted{\unsure{\operatorname{Fr}|\partial}}\partial U$, à fortiori de $\operatorname{Fr}(\deleted{U}X)$, a cette propriété. Soit
% L8–9
\[
\begin{aligned}
\deleted{U}X'=C_{\mathrm{mod}}\unsure{(U)|U U)|U)}=\bigl\{x\in L\setminus U\mid{}&x\text{ est en pos. rel. mod.}\\
&\text{p.r. à }U\bigr\}
\end{aligned}
\]
% L10–13
Alors \deleted{\unsure{$C$|$\mathcal C$}} $\deleted{U}X'$ est \unsure{modérée|modéré}, et c'est la plus grande partie modérée $\deleted{V}Y$ de $L$ telle que $\deleted{V}Y\subset L\setminus\deleted{U}X$, et $Y$ en position relative modérée par. à $X$. Si
% L14
\[
\{T_1\ll T_2<\ldots\ll T_n\}
\]
% L15–20
est le drapeau définissant \deleted{\unsure{$T$|$\Phi$}} $X$, $X'$ est défini par le drapeau formé des $n-1$ tronçons interstitiels \deleted{entre} pour $(T_1,T_2)$, $(T_2,T_3)$ \ldots\ $(T_{n-1},T_n)$, plus (dans la mesure où ils existent) les tronçons, interstitiel gauche de $T_1$, et le tronçon interstitiel droit du $n$. \unsure{Suite|D'autre}
\scanpage{32}
\mspage{31}

% L1–2
Si $n'$ est le cardinal du $ {\unsure{\Phi|\varphi}}'\in\operatorname{Drap}_{\ll}^{\unsure{b|\partial|*}}(\unsure{\operatorname{Tr}|\operatorname{Fr}}(L))$ qui définit $X'$, on aura
% L3
\[
n'\in\{n-1,n,n+1\}
\]
% L4–5
donc
\[
\lvert\operatorname{Card}\unsure{\Pi|\pi}_{\unsure{0|o}}(X)-\operatorname{Card}\Pi_0(X')\rvert\leq 1
\]
% L6–9
De façon précise
\diagram{\(\begin{array}{ll}
\unsure{T|\mathbb T}_1\ \text{non \unsure{majoré|mijoré|minoré}},\ T_n\ \text{non majoré} &: n'=n-1\\
T_1\ \text{minoré et }T_n\ \text{non majoré, ou l'inverse} &: n'=n\\
T_1\ \text{minoré et }T_n\ \text{majoré} &: n'=n+1
\end{array}\)}

% L10–12
On considère l'ens \deleted{ordonné} \deleted{(\unsure{cf. ci-dessous|cf. ci-dessus|p. ci-dessus})} des parties modérées de $L$, $\unsure{\mathcal M|M}=\operatorname{Mod}(L)$. Il y a des structures remarquables

% L13
a) Ordre \quad $X\subset Y$

% L14
b) Relation : $\unsure{\{|\lbrace\!\lbrace}X,Y\unsure{\}|\rbrace\!\rbrace}$ modérée \quad (NB $\unsure{x|X}\subset Y\Rightarrow\{X,Y\}$ \unsure{modéré|modérée})
% Lleft of b–c
% slanted
\margin{Cette relation est-elle déductible de la relation d'ordre -- p.ex. $\{X,Y\}$ mod ssi $X\wedge Y$ existe ???}

% L15–18
c) Partie $\inserted{\mathcal M_f=}\operatorname{Mod}\unsure{\mathrm{fer}|\mathrm{ferm}}(L)$ formée des parties modérées dites ``fermées'' (i.e. $X=X_{\Phi}$, où les $T_i$ $(i\in\Phi)$ sont des \unsure{ps.|ts.} tr. fermés.)

% L19
\uline{Propriétés}
\scanpage{33}
\mspage{32}

% L1–2
1) Pour la relation d'ordre, il y a un plus petit élément $\varnothing$, et un plus grand élément $L$ \inserted{(ou $1$)}.

% L3
2) : $X\subset Y\Rightarrow\{X,Y\}$ modéré

% L4
3) Si $\{X,Y\}$ modéré, $X\wedge Y$ et $X\vee Y$ existent (inf et sup)
% Lleft, opposite 3
% slanted
\margin{Attention $X\vee Y$ \deleted{\unsure{est|n'}} n'est \uline{pas} l'union \deleted{\unsure{de}} $X\cup Y$ !!}

% L5–6
\deleted{\unsure{3)|3 bis)} Si $\{X,\unsure{Y_i|Y}\}$ \unsure{est|sont} \illegible{2} modérés, alors} \deleted{Si \unsure{deux|les} $X$ \unsure{et|sont} \unsure{modérés} \unsure{pour}} \inserted{(distributivité) \unsure{et|si} pour une famille d'\unsure{élmt|élémt} mut\textsuperscript{t} mod, $X\wedge(\unsure{\bigvee_i|\bigvee}Y_i)=\bigvee_i(X\wedge Y_i)$}

% L7–11
\deleted{\unsure{3|5}}4) \deleted{Parmi les} $\forall X\in\unsure{\mathcal M|M}$ \deleted{\unsure{éléments|il existe}}, \deleted{\unsure{dans|du}} parmi les $Y\in\mathcal M$ tels que $\{X,Y\}$ modéré, et $X\wedge Y=\varnothing$, il y a un plus grand élément \inserted{$X'$}. C'est aussi l'unique élément tel que
% Lleft, opposite X′ and tel que
% slanted
\margin{\uline{Dire aussi} Si \deleted{\unsure{(élémt)|($X_i$)}} $\{X_i\}_{i\in\unsure{I|J}}\subset\mathcal M$ est formé d'él\textsuperscript{ts} \unsure{mutuellemt|mutuellement} modérés, alors pour les \unsure{op.} $\wedge,\vee,\unsure{'|{}^c}$ ils engendrent une ``alg. de Boole'', \unsure{finie} si $I$ fini.}
% L12
\[
\{X',X\}\text{ modéré},\qquad \unsure{X'|X''}\wedge X=\varnothing,\qquad X'\vee X=1
\]
% L13–17
Donc l'application $X\mapsto X'$ est involutive, elle renverse la relation d'ordre (anti-involution), \deleted{\unsure{donc} échange} \deleted{\unsure{De|On}}\unsure{De} plus $\{X\ \mathrm{et}\ Y\}$ sont mod. ssi $\{X',Y'\}$ modéré\textsuperscript{*},
% Lleft of same passage
% slanted
\margin{* En fait, $\forall X,Y$ $\{X,Y\}$ mod. ssi $\{X,Y'\}$ mod (donc ssi $\{X',Y'\}$ mod., \unsure{ssi|et} $\{X',Y\}$ modéré)}
% L18
et alors \quad $(X\wedge Y)'=X'\unsure{\cup|\vee}Y'$, \quad $(X\cup Y)'=X'\wedge Y'$

% L19
\deleted{\unsure{4|5}}5) $\varnothing,1\in\mathcal M_{\unsure{f|t}}$

% L20–21
\deleted{5}6) $\forall X\in\mathcal M$, il existe un plus petit $\deleted{\unsure{X|Y}}Y\in\mathcal M_f$ tel que $X\subset Y$ (on pose $Y=\bar X$).
% Lbottom left beside 6
% slanted
\margin{Par \unsure{passage|passant} au compl., on définit $\deleted{\unsure{U|V}}X^{\unsure{\circ|o}}=(\overline{X'})'$.}

% L22–23
De plus Si $X,Y\in\mathcal M$ tels que $\{X,Y\}$ mod, alors $\overline{X\vee Y}=\overline X\vee\overline Y$

% L24 and below
\deleted{6}7) Si $\{X,Y\}$ modéré $\{\overline X,Y\}$ aussi (donc $\{X,\overline Y\}$ et $\{\overline X,\overline Y\}$ aussi
